% ch11_supplement.tex — 第11章：低温工业等离子体物理 补充内容
% ============================================================
% 经典例题 + 完整解答、TikZ图像、核心公式速查卡、自学检查清单

% ============================================================
% 1. 经典例题
% ============================================================

\section{经典例题与解答}

\begin{example}{帕邢曲线分析：最小击穿电压}{ex:paschen}
在平行板电极间充以氩气（Ar），电极间距 $d=2\,\mathrm{cm}$。已知氩气的帕邢常数：$A=12\,\mathrm{cm^{-1}\cdot kPa^{-1}}$，$B=180\,\mathrm{V\cdot cm^{-1}\cdot kPa^{-1}}$，二次电子发射系数 $\gamma=0.1$。求该条件下的最小击穿电压 $V_{min}$ 及对应的优化气压 $p_{min}$。
\end{example}

\textbf{已知条件：}
\begin{itemize}
    \item 电极间距 $d = 2\,\mathrm{cm} = 0.02\,\mathrm{m}$
    \item Townsend第一电离系数经验公式：$\alpha/p = A\exp(-Bp/E)$，即 $\alpha = A p \exp(-Bp/E)$
    \item 帕邢定律：$V_B = \frac{B p d}{\ln(A p d / \ln(1+1/\gamma))}$
\end{itemize}

\textbf{解答：}

帕邢定律给出了击穿电压 $V_B$ 与气压 $p$ 和电极间距 $d$ 乘积 $pd$ 的关系：
\begin{equation}
    V_B = \frac{B \cdot p d}{\ln(A p d) - \ln\left(\ln\left(1+\frac{1}{\gamma}\right)\right)}
\label{eq:ch11-paschen-law}
\end{equation}

为求最小击穿电压，令 $x = pd$。首先计算 $\ln(1+1/\gamma)$：
\[
    \ln\left(1+\frac{1}{0.1}\right) = \ln(11) = 2.398
\]

帕邢定律改写为：
\[
    V_B = \frac{B x}{\ln(A x) - \ln(2.398)} = \frac{B x}{\ln(A x / 2.398)}
\]

对 $V_B$ 关于 $x$ 求导并令其为零。令 $u = A x / 2.398$，则 $V_B = \frac{B x}{\ln u}$。利用 $x = 2.398 u / A$，得：
\[
    V_B = \frac{B \cdot 2.398 u}{A \ln u}
\]

求最小值：令 $f(u) = u / \ln u$，则 $f'(u) = (\ln u - 1) / (\ln u)^2 = 0$，即 $\ln u = 1$，$u = e$。

因此最优条件为：
\[
    \frac{A p d}{\ln(1+1/\gamma)} = e \quad\Rightarrow\quad A p d = e \cdot \ln(1+1/\gamma)
\]

代入数值：
\[
    p d = \frac{e \times 2.398}{A} = \frac{2.718 \times 2.398}{12} = \frac{6.518}{12} = 0.543\,\mathrm{kPa\cdot cm}
\]

因此最优气压（$d=2\,\mathrm{cm}$）：
\[
    p_{min} = \frac{0.543}{d} = \frac{0.543}{2} = 0.272\,\mathrm{kPa}
\]

换算为常用单位：
\[
    \boxed{p_{min} \approx 0.27\,\mathrm{kPa} \approx 2.0\,\mathrm{Torr}}
\]

最小击穿电压：
\[
    V_{min} = \frac{B \cdot p d}{\ln(e)} = B \cdot p d = B \times 0.543
\]
\[
    V_{min} = 180 \times 0.543 = 97.7\,\mathrm{V}
\]

但注意 $B$ 的单位为 $V\cdot cm^{-1}\cdot kPa^{-1}$，$p d$ 单位为 $kPa\cdot cm$，因此：
\[
    \boxed{V_{min} \approx 98\,\mathrm{V}}
\]

\textbf{验证：} 氩气在 $d=2\,\mathrm{cm}$ 时的典型帕邢最小值确实在 $V_{min}\approx 150\sim200\,\mathrm{V}$（不同数据来源有差异）。对于氖气和氮气，帕邢最小值约 $200\sim300\,\mathrm{V}$。此处结果量级合理，说明在给定参数下，最小击穿电压低于 $200\,\mathrm{V}$。

% ------------------------------------------------------------------------

\begin{example}{鞘层厚度计算}{ex:sheath}
某RF感应耦合等离子体（ICP）源中，氩等离子体参数为：
\begin{itemize}
    \item 电子密度 $n_e = 5\times10^{16}\,\mathrm{m^{-3}}$
    \item 电子温度 $T_e = 3\,\mathrm{eV}$
    \item 离子质量 $m_i = 40\,\mathrm{amu}$（氩离子）
    \item 鞘层电压降 $V_s = 20\,\mathrm{V}$（相对于等离子体电位）
\end{itemize}
分别用Bohm判据和Child-Langmuir定律估算鞘层厚度 $s$。
\end{example}

\textbf{解答：}

\textbf{步骤一：Bohm判据——计算离子进入鞘层的速度}

Bohm速度定义为：
\begin{equation}
    v_B = \sqrt{\frac{k_B T_e}{m_i}}
\label{eq:ch11-bohm-criterion-ex}
\end{equation}

其中 $m_i = 40 \times 1.661\times10^{-27}\,\mathrm{kg} = 6.644\times10^{-26}\,\mathrm{kg}$，$k_B T_e = 3 \times 1.602\times10^{-19}\,\mathrm{J} = 4.806\times10^{-19}\,\mathrm{J}$。

\[
    v_B = \sqrt{\frac{4.806\times10^{-19}}{6.644\times10^{-26}}} = \sqrt{7.23\times10^6} = 2.69\times10^3\,\mathrm{m\cdot s^{-1}}
\]

\[
    \boxed{v_B \approx 2.7\times10^3\,\mathrm{m\cdot s^{-1}}}
\]

\textbf{步骤二：Child-Langmuir定律——计算鞘层厚度}

Child-Langmuir定律给出了平面鞘层中电流与鞘层电压的关系。对于离子电流密度 $J_i$：
\begin{equation}
    J_i = \frac{4}{9}\eps_0 \sqrt{\frac{2e}{m_i}} \frac{V_s^{3/2}}{s^2}
\label{eq:ch11-child-langmuir}
\end{equation}

离子电流密度也可由Bohm通量表示：
\[
    J_i = e n_s v_B = 0.6 e n_e v_B
\]

（$n_s \approx 0.6 n_e$ 为鞘层边缘预鞘区密度，通常取此值。）

因此：
\[
    0.6 e n_e v_B = \frac{4}{9}\eps_0 \sqrt{\frac{2e}{m_i}} \frac{V_s^{3/2}}{s^2}
\]

解出 $s^2$：
\[
    s^2 = \frac{4\eps_0}{9 \times 0.6 e n_e v_B} \sqrt{\frac{2e}{m_i}} V_s^{3/2}
\]

将数值代入（使用SI单位）：
\begin{align*}
    \eps_0 &= 8.854\times10^{-12}\,\mathrm{F\cdot m^{-1}} \\
    e &= 1.602\times10^{-19}\,\mathrm{C} \\
    n_e &= 5\times10^{16}\,\mathrm{m^{-3}} \\
    v_B &= 2.69\times10^3\,\mathrm{m\cdot s^{-1}} \\
    V_s &= 20\,\mathrm{V} \\
    m_i &= 6.644\times10^{-26}\,\mathrm{kg}
\end{align*}

计算 $\sqrt{2e/m_i}$：
\[
    \sqrt{\frac{2\times1.602\times10^{-19}}{6.644\times10^{-26}}} = \sqrt{4.82\times10^6} = 2.19\times10^3\,\mathrm{C^{1/2}\cdot kg^{-1/2}}
\]

计算右侧：
\[
    \frac{4\times8.854\times10^{-12}}{9\times0.6\times1.602\times10^{-19}\times5\times10^{16}\times2.69\times10^3} \times 2.19\times10^3 \times (20)^{1.5}
\]

先计算分母：
\[
    9\times0.6\times1.602\times10^{-19}\times5\times10^{16}\times2.69\times10^3 = 5.4\times1.602\times10^{-19}\times5\times10^{16}\times2.69\times10^3
\]
\[
    = 8.65\times10^{-19}\times5\times10^{16}\times2.69\times10^3 = 4.33\times10^{-2}\times2.69\times10^3 = 116.5
\]

分子：
\[
    4\times8.854\times10^{-12} = 3.54\times10^{-11}
\]

\[
    \frac{3.54\times10^{-11}}{116.5} = 3.04\times10^{-13}
\]

\[
    (20)^{1.5} = 20\times\sqrt{20} = 20\times4.47 = 89.4
\]

\[
    s^2 = 3.04\times10^{-13} \times 2.19\times10^3 \times 89.4 = 3.04\times2.19\times89.4\times10^{-10}
\]
\[
    = 595\times10^{-10} = 5.95\times10^{-8}\,\mathrm{m^2}
\]

\[
    s = \sqrt{5.95\times10^{-8}} = 2.44\times10^{-4}\,\mathrm{m}
\]

\[
    \boxed{s \approx 0.24\,\mathrm{mm}}
\]

\textbf{步骤三：Debye长度对比}

Debye长度：
\begin{equation}
    \lambda_D = \sqrt{\frac{\eps_0 k_B T_e}{n_e e^2}} = 743\sqrt{\frac{T_e\,[\mathrm{eV}]}{n_e\,[\mathrm{cm^{-3}}]}}\,\mathrm{cm} = 743\sqrt{\frac{3}{5\times10^{10}}}\,\mathrm{cm}
\label{eq:ch11-debye-length}
\end{equation}

\[
    = 743\times\sqrt{6\times10^{-11}}\,\mathrm{cm} = 743\times7.75\times10^{-6}\,\mathrm{cm} = 5.76\times10^{-3}\,\mathrm{cm}
\]

\[
    \lambda_D \approx 5.8\times10^{-5}\,\mathrm{m} = 58\,\mathrm{\mu m}
\]

鞘层厚度约为Debye长度的 $4\sim5$ 倍，符合物理预期（$s\sim 4\lambda_D$ 对于中等电压鞘层）。

% ------------------------------------------------------------------------

\begin{example}{磁控溅射中的$\vect{E}\times\vect{B}$漂移速度}{ex:exb}
在平面磁控溅射靶面附近，磁场位形近似为：
\begin{itemize}
    \item 靶面法向电场 $E = 500\,\mathrm{V\cdot cm^{-1}} = 5\times10^4\,\mathrm{V\cdot m^{-1}}$
    \item 平行于靶面的磁场分量 $B = 0.03\,\mathrm{T}$
    \item 电子温度 $T_e = 5\,\mathrm{eV}$
\end{itemize}
估算电子的 $\vect{E}\times\vect{B}$ 漂移速度，并与电子热速度比较。
\end{example}

\textbf{解答：}

\textbf{（1）$\vect{E}\times\vect{B}$ 漂移速度}

对于垂直的 $\vect{E}$ 和 $\vect{B}$ 场（靶面法向 $E$ 与靶面切向 $B$ 互相垂直），漂移速度为：
\begin{equation}
    v_{E\times B} = \frac{E}{B}
\label{eq:ch11-exb-drift}
\end{equation}

代入数值：
\[
    v_{E\times B} = \frac{5\times10^4}{0.03} = 1.67\times10^6\,\mathrm{m\cdot s^{-1}}
\]

\[
    \boxed{v_{E\times B} \approx 1.7\times10^6\,\mathrm{m\cdot s^{-1}}}
\]

\textbf{（2）电子热速度}

电子热速度（最概然速率）为：
\begin{equation}
    v_{th,e} = \sqrt{\frac{2 k_B T_e}{m_e}} = 5.93\times10^5 \sqrt{T_e\,[\mathrm{eV}]}
\label{eq:ch11-thermal-velocity}
\end{equation}

\[
    v_{th,e} = 5.93\times10^5 \times \sqrt{5} = 5.93\times10^5 \times 2.236
\]
\[
    = 1.32\times10^6\,\mathrm{m\cdot s^{-1}}
\]

\[
    \boxed{v_{th,e} \approx 1.3\times10^6\,\mathrm{m\cdot s^{-1}}}
\]

\textbf{（3）比较分析}

\[
    \frac{v_{E\times B}}{v_{th,e}} = \frac{1.7\times10^6}{1.3\times10^6} \approx 1.3
\]

$\vect{E}\times\vect{B}$ 漂移速度与电子热速度量级相近。在磁控溅射中，这一漂移使电子被束缚在靶面附近（跑道效应），显著延长电子路径，提高电离效率。当 $v_{E\times B} \sim v_{th,e}$ 时，电子的导向中心运动与无规热运动叠加，形成复杂的漂移轨迹。

% ------------------------------------------------------------------------

\begin{example}{辉光放电等离子体密度估算}{ex:glow}
某直流辉光放电管中：
\begin{itemize}
    \item 放电电压 $V = 500\,\mathrm{V}$
    \item 放电电流 $I = 20\,\mathrm{mA}$
    \item 放电管长度 $L = 30\,\mathrm{cm} = 0.3\,\mathrm{m}$
    \item 管内直径 $D = 2.5\,\mathrm{cm}$，半径 $r = 1.25\,\mathrm{cm} = 0.0125\,\mathrm{m}$
    \item 充入氩气，气压 $p = 100\,\mathrm{Pa}$ ($\approx 0.75\,\mathrm{Torr}$)
\end{itemize}
假设正柱区均匀发光，且电子温度 $T_e \approx 2\,\mathrm{eV}$。估算正柱区平均电子密度 $n_e$。
\end{example}

\textbf{解答：}

\textbf{步骤一：计算放电功率与体积}

放电功率：
\[
    P = I V = 0.02 \times 500 = 10\,\mathrm{W}
\]

正柱区体积（近似为圆柱体）：
\[
    V_p = \pi r^2 L = \pi \times (0.0125)^2 \times 0.3 = \pi \times 1.56\times10^{-4} \times 0.3
\]
\[
    = 1.47\times10^{-4}\,\mathrm{m^3} = 147\,\mathrm{cm^3}
\]

\textbf{步骤二：利用功率平衡估算密度}

在辉光放电中，电子从电场获得的功率主要用于电离和激发。每产生一个电子-离子对，电子需要克服电离能（氩的电离能 $E_{ion} = 15.76\,\mathrm{eV}$）。

若假设放电功率中约 $50\%$ 用于电离（其余用于激发、辐射等），则有效电离功率：
\[
    P_{ion} = 0.5 P = 5\,\mathrm{W}
\]

电离速率（每秒产生的离子对数）：
\[
    R_{ion} = \frac{P_{ion}}{e E_{ion}} = \frac{5}{1.602\times10^{-19} \times 15.76}
\]
\[
    = \frac{5}{2.52\times10^{-18}} = 1.98\times10^{18}\,\mathrm{s^{-1}}
\]

在稳态下，电离速率等于复合速率（在较低气压下以壁面复合为主）：
\[
    R_{ion} = n_e n_i \langle\sigma v\rangle_{rec} V_p + \text{壁面损失}
\]

对于辉光放电正柱区，更简单的估算采用 ambipolar 扩散模型：

正柱区的电子密度可由电流连续性估算。放电电流密度 $J = I / (\pi r^2)$：
\[
    J = \frac{0.02}{\pi \times (0.0125)^2} = \frac{0.02}{4.91\times10^{-4}} = 40.7\,\mathrm{A\cdot m^{-2}}
\]

电子电流密度 $J_e \approx J$（离子电流贡献很小）。
\[
    J_e = e n_e v_{dr,e}
\]

其中电子漂移速度 $v_{dr,e} = -\mu_e E$，电子迁移率 $\mu_e$ 在 $100\,\mathrm{Pa}$ 氩气中约 $0.4\,\mathrm{m^2/(V\cdot s)}$（粗略估计）。

电场强度（正柱区）：
\[
    E_{col} = \frac{V_{col}}{L}
\]

辉光放电中，阴极位降区约占 $200\sim300\,\mathrm{V}$，阳极位降区较小，正柱区电压约 $200\sim300\,\mathrm{V}$。取 $V_{col} = 250\,\mathrm{V}$：
\[
    E_{col} = \frac{250}{0.3} = 833\,\mathrm{V\cdot m^{-1}}
\]

电子漂移速度：
\[
    v_{dr,e} = \mu_e E_{col} = 0.4 \times 833 = 333\,\mathrm{m\cdot s^{-1}}
\]

因此：
\[
    n_e = \frac{J}{e v_{dr,e}} = \frac{40.7}{1.602\times10^{-19} \times 333}
\]
\[
    = \frac{40.7}{5.33\times10^{-17}} = 7.64\times10^{17}\,\mathrm{m^{-3}}
\]

\[
    \boxed{n_e \approx 8\times10^{17}\,\mathrm{m^{-3}}}
\]

\textbf{验证：} 辉光放电的典型电子密度范围为 $10^{15}\sim10^{18}\,\mathrm{m^{-3}}$。对于 $20\,\mathrm{mA}$、$500\,\mathrm{V}$ 的氩气辉光放电，$n_e\sim 10^{17}\,\mathrm{m^{-3}}$ 是合理量级。若气压更高或电流更大，密度会相应增加。

% ============================================================
% 2. TikZ图像
% ============================================================

% ch11 新例题：射频鞘层厚度与自偏压
\begin{example}{射频鞘层与自偏压估算}{rf-sheath-bias}
电容耦合射频放电（CCP）中，极板间距 $d = 3\,\mathrm{cm}$，射频频率 $f = 13.56\,\mathrm{MHz}$，等离子体密度 $n_e = 10^{16}\,\mathrm{m^{-3}}$，电子温度 $T_e = 3\,\mathrm{eV}$。
（1）解释为什么电极会自动形成负自偏压 $V_{dc}$；
（2）估算射频鞘层厚度 $s_m$（设鞘层压降 $V_{rf}\approx 200\,\mathrm{V}$）；
（3）计算鞘层厚度与极板间距之比，说明 CCP 中体等离子区的厚度。
\tcblower
\textbf{解答：}

\textbf{(1) 自偏压的形成：} 电子质量小、响应快，能跟随 $13.56\,\mathrm{MHz}$ 的射频场振荡；离子质量大，只响应时间平均场。电极在射频正半周吸收大量电子（电子通量远大于离子通量），充电为负；负半周电子被排斥，离子缓慢中和。稳态下电极自动积累负电荷，直到\textbf{时间平均电子通量等于离子通量}——这就是\textbf{自偏压}（self-bias）。对两极板面积不等的非对称 CCP，自偏压可达数百伏，方向使小电极（放电功率电极）更负，离子在该电极上获得更高的轰击能量。

\textbf{(2) 射频鞘层厚度：} 用 Child-Langmuir 简化式（$n_s\approx n_e$，$V_0 = 200\,\mathrm{V}$，$T_e = 3\,\mathrm{eV}$，氩 $m_i = 40\,\mathrm{amu}$）：
\[
s \approx \frac{\sqrt{2}}{3}\,\lambda_D\left(\frac{2V_0}{T_e}\right)^{3/4}
\]
德拜长度（$\lambda_D[\mathrm{m}]\approx7430\sqrt{T[\mathrm{eV}]/n[\mathrm{m^{-3}}]}$）：
\[
\lambda_D = 7430\sqrt{\frac{3}{10^{16}}}\,\mathrm{m} = 7430\times5.48\times10^{-8}\,\mathrm{m} \approx 0.41\,\mathrm{mm}
\]
\[
\left(\frac{400}{3}\right)^{3/4} = 133^{0.75} \approx 39
\]
\[
s \approx 0.47\times0.41\,\mathrm{mm}\times39 \approx 7.5\,\mathrm{mm}
\]

\textbf{(3) 鞘层与间距之比：}
\[
\frac{s_m}{d} \approx \frac{7.5\,\mathrm{mm}}{30\,\mathrm{mm}} \approx 0.25 < 1
\]
鞘层明显薄于极板间距（但不完全是"薄鞘"极限），体等离子区仍有约 $15\,\mathrm{mm}$ 的准中性区。若提高密度到 $10^{17}\,\mathrm{m^{-3}}$，$\lambda_D$ 缩小 $\sqrt{10}$ 倍，$s\sim2.4\,\mathrm{mm}$，$s/d\approx0.08$，薄鞘近似良好。这就是工业 CCP 在高密度下加工更均匀的原因：鞘层薄，离子加速区窄，入射能量分散小，各向异性刻蚀更好。
\end{example}

\section{DC辉光放电空间结构图}

\begin{figure}[htbp]
    \centering
    \begin{tikzpicture}[scale=1.0, >=Stealth]
        % 电极
        \draw[thick, fill=gray!30] (-0.5,2) rectangle (0,4);
        \node[rotate=90] at (-0.25,3) {\small 阴极};
        
        \draw[thick, fill=gray!30] (10,2) rectangle (10.5,4);
        \node[rotate=90] at (10.25,3) {\small 阳极};
        
        % 放电区域底色（整体辉光）
        \fill[violet!5] (0,2) rectangle (10,4);
        
        % 各区域（从左到右）
        % 阴极暗区（Aston暗区，非常薄）
        \fill[blue!20] (0,2) rectangle (0.3,4);
        \node[blue] at (0.15,1.5) {\tiny 阴极};
        \node[blue] at (0.15,1.2) {\tiny 暗区};
        
        % 负辉光（最亮）
        \fill[yellow!50] (0.3,2) rectangle (1.5,4);
        \node[orange] at (0.9,1.5) {负辉光区};
        \node[orange] at (0.9,1.1) {\small （最亮）};
        
        % 法拉第暗区
        \fill[gray!15] (1.5,2) rectangle (2.5,4);
        \node[gray] at (2.0,1.5) {\small 法拉第};
        \node[gray] at (2.0,1.2) {\small 暗区};
        
        % 正柱区（均匀发光）
        \fill[green!20] (2.5,2) rectangle (8.5,4);
        \node[green!50!black] at (5.5,3) {\large 正柱区};
        \node[green!50!black] at (5.5,2.5) {\small （均匀等离子体）};
        
        % 阳极暗区（很薄）
        \fill[blue!10] (8.5,2) rectangle (9.2,4);
        \node[blue] at (8.85,1.5) {\small 阳极};
        \node[blue] at (8.85,1.2) {\small 暗区};
        
        % 阳极辉光
        \fill[yellow!30] (9.2,2) rectangle (10,4);
        \node[orange] at (9.6,1.5) {\small 阳极辉光};
        
        % 电位分布示意（上方）
        \draw[thick, red] (0,5.5) -- (0.3,5.3) -- (1.5,4.0) -- (2.5,3.8) -- (8.5,3.2) -- (9.2,3.1) -- (10,3.0);
        \node[red] at (5,5.5) {空间电位分布};
        
        % 电位标注
        \draw[<->, red] (-1,5.5) -- (-1,3.0);
        \node[red, left] at (-1,4.25) {$V_{cathode}$};
        \node[red, below] at (-1,3.0) {$V_{anode}$};
        
        % 各区域发光强度示意（下方柱状图）
        \draw[->] (-0.5,0) -- (10.5,0) node[right] {$x$};
        \draw[->] (-0.5,0) -- (-0.5,1.5) node[above] {\small 发光强度};
        
        % 负辉光峰值
        \fill[orange!50] (0.3,0) rectangle (1.5,1.2);
        \draw[orange, thick] (0.3,0) -- (1.5,1.2) -- (1.5,0);
        
        % 正柱区中等
        \fill[green!30] (2.5,0) rectangle (8.5,0.6);
        \draw[green!50!black, thick] (2.5,0) -- (8.5,0.6) -- (8.5,0);
        
        % 阳极辉光小峰
        \fill[orange!30] (9.2,0) rectangle (10,0.4);
        \draw[orange, thick] (9.2,0) -- (10,0.4) -- (10,0);
        
        % 图例
        \node[anchor=west] at (11,3.5) {\color{blue}\small 暗区：无激发};
        \node[anchor=west] at (11,2.8) {\color{orange}\small 辉光：强激发};
        \node[anchor=west] at (11,2.1) {\color{green!50!black}\small 正柱：等离子体};
        \node[anchor=west] at (11,1.4) {\color{red}\small 曲线：电位分布};
    \end{tikzpicture}
    \caption{DC辉光放电空间结构示意图。从阴极到阳极依次可见：阴极暗区、负辉光区（最亮）、法拉第暗区、正柱区（均匀等离子体）、阳极暗区、阳极辉光。上图红线为空间电位分布，下方柱状图示意发光强度。}
    \label{fig:glow}
\end{figure}

% ============================================================
% 3. 核心公式速查卡
% ============================================================

\section{核心公式速查卡}

\begin{formula}{Townsend放电与帕邢定律}{f:townsend}
Townsend第一电离系数：
\[
    \alpha = A p \exp\left(-\frac{B p}{E}\right)
\]

击穿条件：
\[
    \gamma\left(e^{\alpha d}-1\right) = 1 \quad\Rightarrow\quad \alpha d = \ln\left(1+\frac{1}{\gamma}\right)
\]

帕邢定律（击穿电压）：
\[
    V_B = \frac{B p d}{\ln(A p d) - \ln\left(\ln(1+1/\gamma)\right)}
\]

最小击穿电压条件：
\[
    (pd)_{min} = \frac{e}{A}\ln\left(1+\frac{1}{\gamma}\right), \quad V_{min} = \frac{e B}{A}\ln\left(1+\frac{1}{\gamma}\right)
\]
\end{formula}

\begin{formula}{鞘层理论}{f:sheath}
Bohm判据（鞘层边缘离子速度）：
\[
    v_s \geq v_B = \sqrt{\frac{k_B T_e}{m_i}}
\]

Debye长度：
\begin{equation}
    \lambda_D = \sqrt{\frac{\eps_0 k_B T_e}{n_e e^2}} \approx 743\sqrt{\frac{T_e\,[\mathrm{eV}]}{n_e\,[\mathrm{cm^{-3}}]}}\,\mathrm{cm}
\label{eq:ch11-debye-length-ex}
\end{equation}

Child-Langmuir定律（平面鞘层）：
\begin{equation}
    J_i = \frac{4}{9}\eps_0 \sqrt{\frac{2e}{m_i}} \frac{V_s^{3/2}}{s^2}
\label{eq:ch11-child-langmuir-ex}
\end{equation}

鞘层厚度估算：
\[
    s \approx \lambda_D \left(\frac{2V_s}{T_e}\right)^{3/4} \quad\text{(高电压鞘层)}
\]
\end{formula}

\begin{formula}{工业等离子体中的漂移与扩散}{f:industrial}
$\vect{E}\times\vect{B}$ 漂移速度：
\[
    \vect{v}_{E\times B} = \frac{\vect{E}\times\vect{B}}{B^2}
\]

电子迁移率（与碰撞频率 $\nu$）：
\[
    \mu_e = \frac{e}{m_e \nu}
\]

ambipolar扩散系数：
\[
    D_a = \frac{\mu_i D_e + \mu_e D_i}{\mu_i + \mu_e} \approx D_i\left(1+\frac{T_e}{T_i}\right)
\]
\end{formula}

\begin{formula}{辉光放电正柱区}{f:positive}
电子密度与电流关系：
\[
    J = e n_e v_{dr,e} = e n_e \mu_e E
\]

正柱区电场（维持条件）：
\[
    E \approx \frac{E_{ion}}{\lambda_i \eta_{ion}}
\]
（$\lambda_i$ 为电离平均自由程，$\eta_{ion}$ 为每次电离产生的新电子数）

电子温度与气压关系（经验）：
\[
    k_B T_e \approx \frac{e E_{ion}}{\ln(n_g \sigma_{ion} L)}
\]
\end{formula}

% ============================================================
% 4. 自学检查清单
% ============================================================

\section{自学检查清单}

\begin{checklist}{第11章自学检查清单}{ch11-checklist}
\begin{itemize}[leftmargin=*,label=$\square$]
\item[$\square$] 我能解释Townsend放电的两种电离机制：$\alpha$ 过程和 $\gamma$ 过程
\item[$\square$] 我能画出帕邢曲线，解释为什么存在最小击穿电压
\item[$\square$] 我能根据帕邢定律计算给定 $pd$ 条件下的击穿电压
\item[$\square$] 我知道辉光放电正柱区、负辉光区、阴极暗区的物理成因
\item[$\square$] 我能用Bohm判据估算离子进入鞘层的速度
\item[$\square$] 我能计算Debye长度，并理解鞘层厚度与 $\lambda_D$ 的关系
\item[$\square$] 我能应用Child-Langmuir定律估算鞘层厚度
\item[$\square$] 我了解$\vect{E}\times\vect{B}$漂移在磁控溅射中的作用
\item[$\square$] 我能解释为什么磁控溅射中电子被束缚在靶面附近（跑道效应）
\item[$\square$] 我能从放电功率和电流估算辉光放电的电子密度
\item[$\square$] 我了解ambipolar扩散的物理意义，以及为什么 $D_a \approx D_i(T_e/T_i)$
\item[$\square$] 我能区分电容耦合（CCP）与电感耦合（ICP）等离子体源的基本原理
\end{itemize}
\end{checklist}

\endinput
